% !TEX root = ../main.tex
% Blockchain data used in the thesis: one lane per analysis, one bar per kind of
% on-chain record, placed on the months it covers (1 Oct 2023 to 30 Sep 2026).
% The 27 months from October 2023 to December 2025 are drawn at a compressed scale
% (0.1 per month); the nine months of 2026 at 0.8 per month.
\begin{tikzpicture}[x=1cm, y=1cm, font=\scriptsize, oneline/.style={text height=1.6ex, text depth=0.45ex}]
\definecolor{dsApproval}{RGB}{0,114,178}
\definecolor{dsTransfer}{RGB}{0,158,115}
\definecolor{dsGMX}{RGB}{230,159,0}
\definecolor{dsOther}{RGB}{140,140,140}
\def\xa{1.95}% 1 Oct 2023
\def\xj{4.65}% 1 Jan 2026 (start of the uncompressed scale)
\def\xw{7.05}% 1 Apr 2026: W0 freeze t1 (31 Mar 2026, 23:59:59 UTC)
\def\xt{9.45}% 1 Jul 2026: T_obs and W1 freeze (30 Jun 2026, 23:59:59 UTC)
\def\xe{11.85}% 1 Oct 2026: end of the W1 labels
\def\bh{0.36}% bar height
\def\xr{12.45}% left edge of the right-hand column
% \dsInput{colour}{y}{x start}{x end}{text}: records used as input
% (scores, cohorts, same-window checks).
\def\dsInput#1#2#3#4#5{%
  \fill[#1!35] (#3,{#2-\bh/2}) rectangle (#4,{#2+\bh/2});
  \draw[#1!75!black, line width=0.5pt] (#3,{#2-\bh/2}) rectangle (#4,{#2+\bh/2});
  \node[oneline, inner sep=0pt] at ({(#3+#4)/2},#2) {#5};}
% \dsHeld{colour}{y}{x start}{x end}{text}: labels counted after a freeze date
% and kept out of the scores.
\def\dsHeld#1#2#3#4#5{%
  \fill[pattern=north east lines, pattern color=#1] (#3,{#2-\bh/2}) rectangle (#4,{#2+\bh/2});
  \draw[#1!75!black, line width=0.7pt] (#3,{#2-\bh/2}) rectangle (#4,{#2+\bh/2});
  \node[oneline, fill=white, inner sep=1pt] at ({(#3+#4)/2},#2) {#5};}
\def\rowlabel#1#2#3{\node[oneline, anchor=east, text=#1!75!black] at ({\xa-0.1},#2) {#3};}
\def\lanetitle#1#2{\node[oneline, anchor=west, font=\footnotesize\bfseries, fill=white, inner sep=1pt] at (0,#1) {#2};}
\def\laneuse#1#2{\node[anchor=north west, text width=4.9cm, align=flush left, inner sep=0pt] at (\xr,#1) {#2};}
\def\lanerule#1{\draw[black!25] (0,#1) -- (17.4,#1);}
\def\ytop{-0.3}% top of the lanes
\def\ybot{-9.45}% bottom of the lanes

% Calendar: year boundaries on the compressed part, month boundaries in 2026
\foreach \x in {1.95,2.25,3.45} {\draw[black!12] (\x,\ytop) -- (\x,\ybot);}
\foreach \k in {0,...,9} {\draw[black!12] ({\xj+0.8*\k},\ytop) -- ({\xj+0.8*\k},\ybot);}
\foreach \k/\name in {0/1월,1/2월,2/3월,3/4월,4/5월,5/6월,6/7월,7/8월,8/9월}
  {\node[oneline] at ({\xj+0.8*(\k+0.5)},0) {\name};}
\node[oneline] at ({(\xa+\xj)/2-0.1},0) {줄여 그린 눈금};
\node[oneline, font=\scriptsize\bfseries] at ({(\xa+\xj)/2},0.34) {2023년 10월 -- 2025년 12월};
\node[oneline, font=\scriptsize\bfseries] at ({\xj+0.8*4.5},0.34) {2026년};
\draw[black!60] ({\xj-0.09},-0.24) -- ({\xj-0.01},-0.08);
\draw[black!60] ({\xj+0.01},-0.24) -- ({\xj+0.09},-0.08);
\draw[decorate, decoration={brace, amplitude=4pt}] ({\xa+0.03},0.62) -- ({\xt-0.03},0.62)
  node[oneline, midway, above=4pt] {관측 기간(33개월)};
\draw[decorate, decoration={brace, amplitude=4pt}] ({\xt+0.03},0.62) -- ({\xe-0.03},0.62)
  node[oneline, midway, above=4pt] {W1 라벨 기간};

% Freeze dates
\draw[black!70, dashed, line width=0.6pt] (\xw,\ytop) -- (\xw,\ybot);
\draw[black!70, dashed, line width=0.6pt] (\xt,\ytop) -- (\xt,\ybot);
\node[align=center, anchor=north] at (\xw,{\ybot-0.05}) {2026년 3월 31일\\W0 동결일 $t_1$};
\node[align=center, anchor=north] at (\xt,{\ybot-0.05}) {2026년 6월 30일\\$T_{\mathrm{obs}}$, W1 동결일};

% Lane A: cohort selection
\lanetitle{-0.6}{매칭 코호트 고르기}
\rowlabel{dsApproval}{-1.0}{허락}
\dsInput{dsApproval}{-1.0}{\xa}{\xt}{적어도 세 주인이 해 준 다섯 토큰의 허락}
\laneuse{-0.5}{적어도 세 주인에게서 다섯 토큰의 0이 아닌 허락을 받은 계약 $15{,}107$개예요(\S\ref{sec:data-sources}).}
\lanerule{-1.35}

% Lane B: same-window analysis
\lanetitle{-1.65}{같은 기간 분석, 매칭 코호트($n=15{,}107$)}
\rowlabel{dsApproval}{-2.05}{허락}
\dsInput{dsApproval}{-2.05}{\xa}{\xt}{엔도스랭크; 허락 점검}
\rowlabel{dsTransfer}{-2.49}{송금}
\dsInput{dsTransfer}{-2.49}{\xa}{\xt}{AWP; 송금 점검과 시빌 점검}
\laneuse{-1.55}{두 종류의 화살표를 함께 쓰면 \mbox{C-PR}과 \mbox{S-PR}이 나와요. RQ1, RQ2, RQ3의 벤치마크, 그리고 RQ5의 같은 기간 부분이에요(\S\S\ref{sec:benchmark-results}--\ref{sec:rank-divergence}).}
\lanerule{-2.84}

% Lane C: holdout W0 on spender contracts
\lanetitle{-3.14}{홀드아웃 W0, 스펜더 계약($n=17{,}050$)}
\rowlabel{dsApproval}{-3.54}{허락}
\dsInput{dsApproval}{-3.54}{\xa}{\xw}{점수 입력}
\dsHeld{dsApproval}{-3.54}{\xw}{\xt}{새 허락 쌍}
\rowlabel{dsTransfer}{-3.98}{송금}
\dsInput{dsTransfer}{-3.98}{\xa}{\xw}{점수 입력}
\dsHeld{dsTransfer}{-3.98}{\xw}{\xt}{새 송금자}
\laneuse{-3.04}{$t_1$에 다섯 토큰의 허락을 가지고 있는 스펜더 계약이에요. 점수와 차수 기준선은 $t_1$까지의 이벤트를 써요. RQ4와 RQ5예요(\S\ref{sec:holdout-protocol}, \S\ref{sec:holdout-results}).}
\lanerule{-4.33}

% Lane D: W0 trader run
\lanetitle{-4.63}{W0 트레이더 실행, 탐색적($n=88$)}
\rowlabel{dsApproval}{-5.03}{허락}
\dsInput{dsApproval}{-5.03}{\xa}{\xw}{점수 입력}
\rowlabel{dsTransfer}{-5.47}{송금}
\dsInput{dsTransfer}{-5.47}{\xa}{\xw}{점수 입력}
\rowlabel{dsGMX}{-5.91}{GMX V2}
\dsHeld{dsGMX}{-5.91}{\xw}{\xt}{트레이더 라벨}
\laneuse{-4.53}{$t_1$ 그래프에 있고, 4월부터 6월까지 포지션을 적어도 세 번 닫은 GMX~V2 계약 계정이에요. 이익 라벨에는 어떤 규칙도 쓰지 않아요(\S\ref{sub:trader-holdout}). 이 계약들의 청산 라벨은 결과 비교에만 들어가요.}
\lanerule{-6.26}

% Lane E: registered window W1
\lanetitle{-6.56}{등록 기간 W1}
\rowlabel{dsApproval}{-6.96}{허락}
\dsInput{dsApproval}{-6.96}{\xa}{\xt}{점수 입력}
\dsHeld{dsApproval}{-6.96}{\xt}{\xe}{새 허락 쌍}
\rowlabel{dsTransfer}{-7.40}{송금}
\dsInput{dsTransfer}{-7.40}{\xa}{\xt}{점수 입력}
\dsHeld{dsTransfer}{-7.40}{\xt}{\xe}{새 송금자}
\rowlabel{dsGMX}{-7.84}{GMX V2}
\dsHeld{dsGMX}{-7.84}{\xt}{\xe}{청산 없이 닫기}
\laneuse{-6.46}{스펜더 $18{,}152$개와 트레이더 $70$개이고, 모두 계약이에요. 고친 점수로 규칙~B와 결과 비교를 해요. 이 연구가 2026년 10월 9일 16:49~UTC의 등록 뒤에 뽑아냈어요(\S\ref{sub:fresh-holdout}).}
\lanerule{-8.19}

% Lane F: account type
\lanetitle{-8.49}{계정 종류}
\rowlabel{dsOther}{-8.95}{코드}
\fill[dsOther!75!black] (12.1,-8.95) +(0,0.13) -- +(0.13,0) -- +(0,-0.13) -- +(-0.13,0) -- cycle;
\node[anchor=east, fill=white, inner sep=1pt] at (11.9,-8.95) {고정한 블록에서 \texttt{eth\_getCode}, 2026년 10월 9일};
\laneuse{-8.39}{후보, 스펜더, 트레이더 주소 가운데 어느 것이 계약인지 알아봐요. EIP-7702 위임 표시는 EOA로 세어요(\S\ref{sec:data-sources}).}

% Legend: record colours, then bar styles
\begin{scope}[shift={(0,-10.55)}]
  \node[oneline, anchor=west, font=\scriptsize\bfseries, inner sep=0pt] at (0,0) {기록};
  \foreach \c/\t/\x/\y in {dsApproval/ERC-20 허락 로그/1.4/0, dsTransfer/ERC-20 송금 로그/5.0/0, dsGMX/GMX V2 포지션 닫기/8.6/0, dsOther/RPC로 읽은 계약 코드/12.45/0} {
    \fill[\c!35] (\x,{\y-0.1}) rectangle ({\x+0.45},{\y+0.1});
    \draw[\c!75!black, line width=0.5pt] (\x,{\y-0.1}) rectangle ({\x+0.45},{\y+0.1});
    \node[oneline, anchor=west] at ({\x+0.5},\y) {\t};
  }
  \node[oneline, anchor=west, font=\scriptsize\bfseries, inner sep=0pt] at (0,-0.5) {막대};
  \fill[dsApproval!35] (1.4,-0.6) rectangle (1.625,-0.4);
  \fill[dsTransfer!35] (1.625,-0.6) rectangle (1.85,-0.4);
  \draw[black!60, line width=0.5pt] (1.4,-0.6) rectangle (1.85,-0.4);
  \node[oneline, anchor=west] at (1.9,-0.5) {칠한 막대: 점수, 코호트, 같은 기간 점검에 넣는 입력};
  \fill[pattern=north east lines, pattern color=black!70] (9.4,-0.6) rectangle (9.85,-0.4);
  \draw[black!60, line width=0.7pt] (9.4,-0.6) rectangle (9.85,-0.4);
  \node[oneline, anchor=west] at (9.9,-0.5) {빗금 막대: 떼어 둔 라벨, 동결일 뒤에 셈};
  \draw[black!70, dashed, line width=0.6pt] (1.625,-1.1) -- (1.625,-0.74);
  \node[oneline, anchor=west] at (1.9,-0.92) {점선: 동결일};
\end{scope}
\end{tikzpicture}
